% Part: first-order-logic
% Chapter: syntax-and-semantics
% Section: main-operator

\documentclass[../../../include/open-logic-section]{subfiles}

\begin{document}

\olfileid[fr]{fol}{syn}{mai}

\olsection{\printtoken{S}{main operator} d'une formule}

\begin{explain}
Il est souvent utile de parler du dernier opérateur employé pour
construire !!a{formula}~$!A$. On l'appelle l'\emph{!!{main operator}}
de~$!A$. Intuitivement, c'est l'opérateur le plus « extérieur »
de~$!A$. Par exemple, l'!!{main operator} de $\lnot !A$ est
$\lnot$, celui de $(!A \lor !B)$ est~$\lor$, et ainsi de suite.
\end{explain}

\begin{defn}[!!^{main operator}]
\ollabel{def:main-op}
L'\emph{!!{main operator}} d'!!a{formula}~$!A$ est défini comme suit :
\begin{enumerate}
\item \indcase*{!A}{!A}{$\indfrm$ n'a pas d'!!{main operator}.}

\tagitem{prvNot}{\indcase{!A}{\lnot !B}{l'!!{main operator} de
  $\indfrm$ est~$\lnot$.}}{}

\tagitem{prvAnd}{\indcase{!A}{(!B \land !C)}{l'!!{main operator} de
  $\indfrm$ est~$\land$.}}{}

\tagitem{prvOr}{\indcase{!A}{(!B \lor !C)}{l'!!{main operator} de
  $\indfrm$ est~$\lor$.}}{}

\tagitem{prvIf}{\indcase{!A}{(!B \lif !C)}{l'!!{main operator} de
  $\indfrm$ est~$\lif$.}}{}

\tagitem{prvIff}{\indcase{!A}{(!B \liff !C)}{l'!!{main operator} de
  $\indfrm$ est~$\liff$.}}{}

\tagitem{prvAll}{\indcase{!A}{\lforall[x][!B]}{l'!!{main operator}
  de $\indfrm$ est~$\lforall$.}}{}

\tagitem{prvEx}{\indcase{!A}{\lexists[x][!B]}{l'!!{main operator}
  de $\indfrm$ est~$\lexists$.}}{}
\end{enumerate}
\end{defn}

Dans chaque cas, il s'agit de l'\emph{occurrence} précisément indiquée
de l'!!{main operator} dans la !!{formula}. Ainsi,
$((!D \lif !E) \lif (!E \lif !D))$ est de la forme $(!B \lif !C)$,
avec $!B$ égal à~$(!D \lif !E)$ et $!C$ égal à~$(!E \lif !D)$ ;
son !!{main operator} est donc la deuxième occurrence de~$\lif$.

\begin{explain}
Il s'agit d'une définition \emph{par cas} d'une fonction qui associe
à chaque !!{formula} non atomique l'occurrence de son
!!{main operator}.\footnote{Précision éditoriale : la source dit « récursive ». Les clauses ci-dessus sélectionnent directement un opérateur sans faire appel à la valeur de cette fonction sur les sous-formules. Il s'agit d'une définition par cas, selon la distinction explicitée dans la section suivante.} La définition inductive des !!{formula}s garantit
que toute !!{formula}~$!A$ relève de l'un des cas de la
\olref{def:main-op} : toute !!{formula} non atomique~$!A$ possède donc
un !!{main operator}. Chaque !!{formula} ne relevant que d'un seul
cas, et les !!{formula}s plus petites dont $!A$ est construite étant
déterminées de manière unique, l'occurrence de l'!!{main operator}
de~$!A$ est elle-même unique. Nous avons donc bien défini une fonction.
\end{explain}

Nous nommons les !!{formula}s selon leur !!{main operator}, comme
l'indique \olref{tab:main-op}.\iftag{defNot,defOr,defAnd,defIf,defIff,defTrue,defFalse,defEx,defAll}
{ Rappelons cependant que les opérateurs définis ne figurent pas
officiellement dans les !!{formula}s : ce sont des abréviations. Ils
ne peuvent donc pas, à proprement parler, être l'!!{main operator}
d'une !!{formula}, et il n'est pas nécessaire de les traiter
séparément dans les démonstrations concernant toutes les
!!{formula}s.}

\begin{table}[!h]
\centering
\begin{tabular}{c | c | c}
!!^{main operator} & Type de !!{formula} & Exemple\\
\hline
aucun & atomique (!!{formula}) &
\iftag{prvFalse}{$\lfalse$,}{}
\iftag{prvTrue}{$\ltrue$,}{}
$\Atom{R}{t_1, \dots, t_n}$,
$\eq[t_1][t_2]$\\
$\lnot$ & négation & $\lnot !A$ \\
$\land$ & conjonction & $(!A \land !B)$ \\
$\lor$ & disjonction & $(!A \lor !B)$ \\
$\lif$ & !!{conditional} & $(!A \lif !B)$ \\
$\liff$ & !!{biconditional} & $(!A \liff !B)$ \\
$\lforall[][]$ & universelle (!!{formula}) & $\lforall[x][!A]$ \\
$\lexists[][]$ & existentielle (!!{formula}) & $\lexists[x][!A]$
\end{tabular}
\caption{Opérateur principal et noms des !!{formula}s}
\ollabel{tab:main-op}
\end{table}

\end{document}
